% 原书第171—174页：导言及9.1。
\chapter{潜变量}
\label{ch:hidden-variables}
\newcounter{chninetransnote}
\providecommand{\chninenote}[1]{\stepcounter{chninetransnote}\begingroup
\renewcommand{\thefootnote}{译注\arabic{chninetransnote}}\footnote{#1}\addtocounter{footnote}{-1}\endgroup}

到目前为止，我们一直假定，模型中的所有变量，除噪声外，都已被测量。实际中，随机变量集合是由我们自己选择的，因此，需要界定哪些变量在“因果上相关”。为此，第9.1节引入“因果充分性”与“干预充分性”两个概念。即使暂且不讨论精确定义的细节，也显然可以看到：在大多数实际应用中，许多因果上相关的变量并未被观测到。这些未观测变量称为\textbf{潜变量（hidden variables，亦称latent variables）}。

Simpson悖论，第9.2节，说明忽略潜在混杂如何导致错误的因果结论。在线性设定下，一种常称为“工具变量”的结构，可以使对应于因果效应的回归系数具有可识别性，见例6.42和第9.3节。如何为含潜变量的结构因果模型找到良好的图表示，尤其是能够编码条件独立结构的表示，仍是一个活跃的研究领域；第9.4节将介绍其中一些方案。最后，潜变量还会使观测分布中出现超出条件独立关系的约束，见第9.5节。我们将简要讨论如何将这些约束用于结构学习，但不涉及方法实现细节。关于潜变量处理的更多历史说明，参见 Spirtes et al.［2000，第6.1节］。

\section{干预充分性}

通常，如果不存在一个未观测共同原因 $C\notin\mathbf X$，同时导致 $\mathbf X$ 中的多个变量，就称变量集合 $\mathbf X$ 满足\textbf{因果充分性（causal sufficiency）}［例如 Spirtes, 2010］。这一定义符合我们对“相关”变量集合的直观理解，但它使用了“共同原因”概念，因此，应当相对于某个更大的变量集合 $\widetilde{\mathbf X}\supseteq\mathbf X$ 来理解；当然，对于这个更大集合，也可能还需要定义因果充分性。在更大集合 $\widetilde{\mathbf X}$ 所对应的结构因果模型中，如果存在从 $C$ 到 $X$、从 $C$ 到 $Y$ 的有向路径，而且前者不包含 $Y$，后者不包含 $X$，那么 $C$ 就是 $X,Y$ 的共同原因。共同原因也称为混杂因子，以下将交替使用这两个说法。

我们对因果充分性作一个小修改，称之为\textbf{干预充分性（interventional sufficiency）}。这个概念以结构因果模型的可否证性为基础，见第6.8节。

\begin{definition}[干预充分性]
如果存在一个定义在变量集合 $\mathbf X$ 上的结构因果模型，使其作为干预模型不能被否证，就称 $\mathbf X$ 满足干预充分性。也就是说，这个模型诱导的观测分布和干预分布，与实际观测到的分布一致。
\end{definition}

我们认为，这个概念在直观上很有吸引力，因为它描述了一个变量集合何时已经足够大，能够进行因果推理，即计算观测分布和干预分布。

如果存在一个潜在共同原因，那么只考虑两个变量通常不够，这应当不难理解。根据定义，这两个变量不满足因果充分性；而第9.2节的Simpson悖论，另见例6.37，则表明它们一般也不满足干预充分性。事实上，这一悖论把问题推到了极端：在两个观测变量上忽略混杂的结构因果模型，不仅导出错误的干预分布，甚至会颠倒因果效应的正负号，使有害的治疗看起来有益，见式 (9.2)。

不过，有时即使存在潜在混杂，仍然可以计算出正确的干预分布。下面例子中的变量集合，满足干预充分性，却不满足因果充分性。

\begin{example}
考虑以下结构因果模型：
\[
 \begin{aligned}
 Z&:=N_Z,\\
 X&:=\mathbf1_{\{Z\geq2\}}+N_X,\\
 Y&:=Z\bmod2+X+N_Y,
 \end{aligned}
\]
其中，$N_Z\sim U(\{0,1,2,3\})$，即在 $\{0,1,2,3\}$ 上均匀分布，且 $N_X,N_Y\overset{\mathrm{i.i.d.}}{\sim}\mathcal N(0,1)$，见图9.1左图。虽然 $X,Y$ 显然不满足因果充分性，\footnote{这里，潜在共同原因 $Z$ 不仅有箭头指向 $X,Y$，而且对二者都具有总因果效应，见定义6.12。}
但可以证明，这两个变量满足干预充分性。原因在于，这个“混杂因子”$Z$ 由两个独立部分构成：$Z_1:=\mathbf1_{\{Z\geq2\}}$ 是 $Z$ 二进制表示的第一位，$Z_2:=Z\bmod2$ 是第二位。在这个意义上，可以将“混杂因子”分成相互独立的变量 $Z_1,Z_2$，其中 $Z_1$ 影响 $X$，$Z_2$ 影响 $Y$，见图9.1。
\end{example}

\begin{figure}[htbp]
\centering
\begin{tikzpicture}[causal,x=1.6cm,y=1.4cm]
 \node[vtx] (ch09f1az) at (1,1.2) {$Z$};
 \node[vtx] (ch09f1ax) at (0,0) {$X$};
 \node[vtx] (ch09f1ay) at (2,0) {$Y$};
 \draw[ed] (ch09f1az)--(ch09f1ax);
 \draw[ed] (ch09f1az)--(ch09f1ay);
 \draw[ed] (ch09f1ax)--(ch09f1ay);
 \node[vtx] (ch09f1bz1) at (4.0,1.2) {$Z_1$};
 \node[vtx] (ch09f1bz2) at (5.0,1.2) {$Z_2$};
 \node[vtx] (ch09f1bx) at (3.5,0) {$X$};
 \node[vtx] (ch09f1by) at (5.5,0) {$Y$};
 \draw[ed] (ch09f1bz1)--(ch09f1bx);
 \draw[ed] (ch09f1bz2)--(ch09f1by);
 \draw[ed] (ch09f1bx)--(ch09f1by);
\end{tikzpicture}
\caption{对于例9.2的模型，两个图表示干预等价的结构因果模型。虽然只有第二种表示使 $X,Y$ 满足因果充分性，但无论采用哪种表示，$X,Y$ 都满足干预充分性。}
\label{fig:ch09-1}
\end{figure}


一般而言，因果充分性与干预充分性有如下关系，证明见附录C.13。

\begin{proposition}[干预充分性与因果充分性]
设 $\mathfrak C$ 是变量集合 $\mathbf X$ 上的结构因果模型，并且作为干预模型不能被否证。
\begin{enumerate}[label=(\roman*)]
 \item 如果子集 $\mathbf O\subseteq\mathbf X$ 满足因果充分性，那么它也满足干预充分性。
 \item 一般而言，反过来不成立。也就是说，存在满足干预充分性、但不满足因果充分性的集合 $\mathbf O\subseteq\mathbf X$。
\end{enumerate}
\end{proposition}

此外，例9.2表明，判断一个变量子集是否满足干预充分性，不可能仅依靠图判据。对于许多结构类似于图9.1左图的结构因果模型，$X,Y$ 并不满足干预充分性。不过，下面的注说明，省略一个“中间”变量会保留干预充分性。

\begin{remark}
有以下三条结论。
\begin{enumerate}[label=(\roman*)]
 \item 假设存在一个 $X,Y,Z$ 上的结构因果模型，图为 $X\to Y\to Z$，满足 $X\nindep Z$，并且诱导正确的干预分布。那么，$X,Z$ 满足干预充分性，因为可以在 $X,Z$ 上构造图为 $X\to Z$ 的结构因果模型。
 \item 假设存在一个 $X,Y,Z$ 上的结构因果模型 $\mathfrak C$，能够诱导正确的干预分布，图为 $X\to Y\to Z$，并另有边 $X\to Z$；再假设 $P^{\mathfrak C}_{X,Y,Z}$ 忠实于这一图，另见（iii）。那么，同样可以利用 $X,Z$ 上图为 $X\to Z$ 的结构因果模型，说明 $X,Z$ 满足干预充分性。
 \item 假设情形与（ii）相同，但区别在于，对所有 $x$，都有
 \[
  P^{\mathfrak C}_{Z\mid X=x}
  =P_Z^{\mathfrak C;\mathrm{do}(X:=x)}
  =P_Z^{\mathfrak C}.
 \]
 特别地，$P^{\mathfrak C}_{X,Y,Z}$ 不忠实于该图。那么，$X,Z$ 满足干预充分性，因为可以采用 $X,Z$ 上具有空图的结构因果模型。注意，反事实未必能够得到正确表示。
\end{enumerate}
这些结论的证明留给读者，见习题9.10。
\end{remark}

当我们找到一个仅包含观测变量的结构因果模型，并且它与包含全部变量的原始模型干预等价时，可以把前者称为\textbf{边缘化结构因果模型（marginalized SCM）}。前面已经看到，确定边缘化模型的结构，不存在仅依赖图的判据，还需要关于因果机制的信息，也就是结构赋值的具体形式。

Bongers et al.［2016］更详细地研究了结构因果模型的边缘化。核心思路是从原始模型出发，只考虑观测变量的结构赋值；只要赋值右端出现潜变量，就反复代入相应潜变量的结构赋值。这样会得到噪声变量为多元、且可能相依的结构因果模型。在某些情况下，随后可以选择一个干预等价、噪声变量为一元的结构因果模型。
